\chapter{Chain Contractions}
\label{ch:viii23-chain-contractions}

A chain contraction is a chain homotopy from the identity to zero.  It gives an
explicit algebraic certificate that every cycle bounds.  Thus contractible
chain complexes are acyclic, while the converse requires additional splitting
hypotheses.

\section{Definition}
\begin{definition}[Chain contraction]
\label{def:viii23-contraction}
A chain contraction on \(C_\bullet\) is a family
\[
s_n:C_n\to C_{n+1}
\]
such that
\[
\operatorname{id}_C=\partial s+s\partial.
\]
\end{definition}

A complex admitting such an \(s\) is contractible as a chain complex.

\section{Contractible implies acyclic}
\begin{theorem}
\label{thm:viii23-acyclic}
If \(C_\bullet\) admits a chain contraction, then
\[
H_n(C)=0
\]
for all \(n\).
\end{theorem}

\begin{proof}
For a cycle \(z\),
\[
z=\partial s(z)+s(\partial z)=\partial s(z),
\]
so every cycle is a boundary.
\end{proof}

\section{Acyclic versus contractible}
Acyclic means only \(H_*(C)=0\).  Contractible means the stronger identity map
is null-homotopic.  In general:
\[
\text{contractible}\Longrightarrow\text{acyclic},
\]
but not conversely without hypotheses.

\section{Split exact complexes}
\begin{theorem}[Split exact complexes contract]
\label{thm:viii23-split}
An exact chain complex of modules whose short exact sequences
\[
0\to Z_n\to C_n\to B_{n-1}\to0
\]
split admits a chain contraction.
\end{theorem}

\section{Free complexes}
Over a field, every short exact sequence of vector spaces splits.  Hence every
acyclic chain complex of vector spaces is contractible.

For free abelian complexes, splitting must be checked more carefully.

\section{Elementary contractible pair}
Consider
\[
0\to A\xrightarrow{\operatorname{id}}A\to0
\]
in adjacent degrees.  It is contractible.

\section{Disk-like complexes}
The augmented simplicial chain complex of a simplex is contractible.  An
explicit contraction can be obtained by coning every simplex to a fixed vertex.

\section{Cone-to-a-vertex operator}
Fix a vertex \(v\).  On a simplex \([v_0,\dots,v_n]\), define informally
\[
s([v_0,\dots,v_n])=[v,v_0,\dots,v_n]
\]
when this simplex is valid, with the standard conventions when \(v\) is already
present.  Then
\[
\partial s+s\partial=\operatorname{id}
\]
on the reduced chain complex of a simplex.

\section{Reduced versus unreduced contraction}
A nonempty contractible space has ordinary \(H_0\cong\mathbb Z\), so its
ordinary chain complex is not acyclic.  The reduced or augmented chain complex
is the correct object to contract.

\section{Strong deformation retract analogy}
Topologically, a contraction deforms a space to a point.  Algebraically, a
chain contraction deforms the identity chain map to zero.  For reduced chains
of a contractible space, these ideas align.

\section{Chain contraction and null maps}
If \(C\) is contractible, every chain map
\[
f:C\to D
\]
is null-homotopic:
\[
f=f\operatorname{id}_C
=f(\partial s+s\partial)
=\partial(fs)+(fs)\partial.
\]

Similarly, every chain map \(D\to C\) is null-homotopic.

\section{Direct sums}
A direct sum of contractible complexes is contractible by taking the direct sum
of their contractions.

\section{Tensor preview}
Tensoring a contraction with the identity often produces a contraction on a
tensor product, with the usual grading signs.  This becomes useful in VIII/27.

\section{Filtration-compatible contractions}
If a contraction respects a filtration, it can simplify spectral or filtered
computations by collapsing contractible pieces without altering homology.

\section{Cancellation of a basis pair}
In a finite free chain complex, if a boundary matrix contains a unit entry that
pairs one basis generator with another, an elementary chain contraction can
often cancel that pair.

This is the algebraic basis of many chain-complex reduction algorithms.

\section{Gaussian elimination over a field}
Over a field, any nonzero pivot is a unit.  Acyclic finite-dimensional chain
complexes can therefore be reduced to direct sums of elementary contractible
two-term complexes.

\section{Smith normal form over \(\mathbb Z\)}
Over \(\mathbb Z\), diagonal entries equal to \(\pm1\) correspond to
contractible summands, while entries of absolute value greater than \(1\)
produce torsion homology and therefore cannot be contracted away.

\section{Acyclic finite free example}
The complex
\[
0\to\mathbb Z
\xrightarrow{\binom{1}{0}}
\mathbb Z^2
\xrightarrow{(0\ \ 1)}
\mathbb Z
\to0
\]
is exact and splits, hence contractible.

\section{Noncontractible torsion example}
The two-term complex
\[
0\to\mathbb Z\xrightarrow{\times m}\mathbb Z\to0
\]
with \(|m|>1\) has
\[
H_0\cong\mathbb Z/m,
\]
so it is not contractible.

\section{Acyclicity certificate}
A chain contraction is stronger than merely knowing homology vanishes: it gives
a formula that explicitly fills each cycle.

\section{Homological perturbation preview}
If a differential is modified slightly, one may try to correct an existing
contraction rather than recompute homology from scratch.  This idea belongs to
more advanced homological perturbation theory.

\section{Bridge to mapping cones}
Mapping cones detect chain-homotopy equivalences.  In particular, the cone of a
chain map is acyclic when the map is a quasi-isomorphism, and is contractible
when the map is a chain-homotopy equivalence under the standard algebraic
hypotheses.  VIII/24 develops this construction.


\section{Exercises}

\begin{exercise}\label{exr:viii23-01}
Exercise 1
\end{exercise}
\begin{hint}
% VIII-PEDAGOGY-HINT-v1.1 exr:viii23-01 replace_template
Check \(s_n:C_n\to C_{n+1}\) and verify \(\partial_{n+1}s_n+s_{n-1}\partial_n=\operatorname{id}_{C_n}\) in each degree.
\end{hint}
\begin{solution}
A chain contraction satisfies \(\operatorname{id}=\partial s+s\partial\).
\end{solution}

\begin{exercise}\label{exr:viii23-02}
Exercise 2
\end{exercise}
\begin{hint}
% VIII-PEDAGOGY-HINT-v1.1 exr:viii23-02 replace_template
Evaluate the contraction equation on a cycle; the resulting filling shows that every cycle represents the zero homology class.
\end{hint}
\begin{solution}
A contractible chain complex has zero homology.
\end{solution}

\begin{exercise}\label{exr:viii23-03}
Exercise 3
\end{exercise}
\begin{hint}
% VIII-PEDAGOGY-HINT-v1.1 exr:viii23-03 replace_template
When \(\partial z=0\), compute \(\partial s(z)\); the operator provides an explicit filling rather than just an existence argument.
\end{hint}
\begin{solution}
For a cycle \(z\), the formula gives \(z=\partial s(z)\).
\end{solution}

\begin{exercise}\label{exr:viii23-04}
Exercise 4
\end{exercise}
\begin{hint}
% VIII-PEDAGOGY-HINT-v1.1 exr:viii23-04 replace_template
Test the exact complex \(0\to\mathbb Z\xrightarrow{2}\mathbb Z\to\mathbb Z/2\to0\); a contraction would force a splitting that does not exist.
\end{hint}
\begin{solution}
Acyclicity alone does not always provide a contraction.
\end{solution}

\begin{exercise}\label{exr:viii23-05}
Exercise 5
\end{exercise}
\begin{hint}
% VIII-PEDAGOGY-HINT-v1.1 exr:viii23-05 replace_template
Write \(C_n=B_n\oplus L_n\) with \(\partial:L_n\to B_{n-1}\) an isomorphism, and use its inverse on the boundary summand to define \(s\).
\end{hint}
\begin{solution}
Split exact complexes admit contractions.
\end{solution}

\begin{exercise}\label{exr:viii23-06}
Exercise 6
\end{exercise}
\begin{hint}
% VIII-PEDAGOGY-HINT-v1.1 exr:viii23-06 replace_template
Choose a complement to each cycle subspace. Exactness identifies the differential on that complement with the boundary subspace one degree lower.
\end{hint}
\begin{solution}
Every acyclic complex of vector spaces is contractible.
\end{solution}

\begin{exercise}\label{exr:viii23-07}
Exercise 7
\end{exercise}
\begin{hint}
% VIII-PEDAGOGY-HINT-v1.1 exr:viii23-07 replace_template
Send the lower copy of \(A\) back to the upper copy by the identity, set other components to zero, and verify the equation on both copies.
\end{hint}
\begin{solution}
The elementary complex \(0\to A\xrightarrow{\mathrm{id}}A\to0\) is contractible.
\end{solution}

\begin{exercise}\label{exr:viii23-08}
Exercise 8
\end{exercise}
\begin{hint}
% VIII-PEDAGOGY-HINT-v1.1 exr:viii23-08 replace_template
Choose a cone vertex and prepend it to an oriented simplex, with zero assigned to repeated-vertex simplices; include the augmented degree in the boundary check.
\end{hint}
\begin{solution}
The reduced chain complex of a simplex is contractible.
\end{solution}

\begin{exercise}\label{exr:viii23-09}
Exercise 9
\end{exercise}
\begin{hint}
% VIII-PEDAGOGY-HINT-v1.1 exr:viii23-09 replace_template
Ordinary degree-zero homology of a point is nonzero, so its identity cannot be null-homotopic; augmentation removes this obstruction.
\end{hint}
\begin{solution}
Ordinary chains of a point retain \(H_0\), so reduced chains are used.
\end{solution}

\begin{exercise}\label{exr:viii23-10}
Exercise 10
\end{exercise}
\begin{hint}
% VIII-PEDAGOGY-HINT-v1.1 exr:viii23-10 replace_template
For a chain map \(f:C\to D\), test \(fs\) as the null-homotopy operator and use that \(f\) commutes with differentials.
\end{hint}
\begin{solution}
If \(C\) is contractible, every map \(C\to D\) is null-homotopic.
\end{solution}

\begin{exercise}\label{exr:viii23-11}
Exercise 11
\end{exercise}
\begin{hint}
% VIII-PEDAGOGY-HINT-v1.1 exr:viii23-11 replace_template
For a chain map \(f:D\to C\), test \(sf\), using the contraction on the target rather than on the source.
\end{hint}
\begin{solution}
If \(C\) is contractible, every map \(D\to C\) is null-homotopic.
\end{solution}

\begin{exercise}\label{exr:viii23-12}
Exercise 12
\end{exercise}
\begin{hint}
% VIII-PEDAGOGY-HINT-v1.1 exr:viii23-12 replace_template
Take the direct sum of the degree-raising operators and check the identity separately on every summand, retaining finite support.
\end{hint}
\begin{solution}
Direct sums of contractions remain contractions.
\end{solution}

\begin{exercise}\label{exr:viii23-13}
Exercise 13
\end{exercise}
\begin{hint}
% VIII-PEDAGOGY-HINT-v1.1 exr:viii23-13 replace_template
First isolate a unit boundary coefficient by compatible basis operations; then invert that unit to contract the paired generators and adjust neighboring maps.
\end{hint}
\begin{solution}
A unit pivot in a boundary matrix often determines a cancellable pair.
\end{solution}

\begin{exercise}\label{exr:viii23-14}
Exercise 14
\end{exercise}
\begin{hint}
% VIII-PEDAGOGY-HINT-v1.1 exr:viii23-14 replace_template
A nonzero field element has an inverse, whereas an integer pivot such as two cannot be inverted over \(\mathbb Z\).
\end{hint}
\begin{solution}
Over a field every nonzero pivot is a unit.
\end{solution}

\begin{exercise}\label{exr:viii23-15}
Exercise 15
\end{exercise}
\begin{hint}
% VIII-PEDAGOGY-HINT-v1.1 exr:viii23-15 replace_template
On an isolated two-term block with differential \(\pm1\), its inverse supplies the degree-raising contraction of that block.
\end{hint}
\begin{solution}
Smith entries \(\pm1\) correspond to contractible summands.
\end{solution}

\begin{exercise}\label{exr:viii23-16}
Exercise 16
\end{exercise}
\begin{hint}
% VIII-PEDAGOGY-HINT-v1.1 exr:viii23-16 replace_template
For an isolated block \(\mathbb Z\xrightarrow{m}\mathbb Z\), compute the cokernel; a nonzero torsion class rules out a contraction.
\end{hint}
\begin{solution}
Smith entries \(|m|>1\) can produce torsion homology.
\end{solution}

\begin{exercise}\label{exr:viii23-17}
Exercise 17
\end{exercise}
\begin{hint}
% VIII-PEDAGOGY-HINT-v1.1 exr:viii23-17 replace_template
Given a cycle \(z\), the certificate must produce the chain \(s(z)\) and verify its boundary is exactly \(z\).
\end{hint}
\begin{solution}
A contraction is an explicit certificate that cycles bound.
\end{solution}

\begin{exercise}\label{exr:viii23-18}
Exercise 18
\end{exercise}
\begin{hint}
% VIII-PEDAGOGY-HINT-v1.1 exr:viii23-18 replace_template
Since \(\partial\) lowers degree, \(s\) must raise it for the two composites in the contraction equation to act in the original degree.
\end{hint}
\begin{solution}
A contraction has degree \(+1\).
\end{solution}

\begin{exercise}\label{exr:viii23-19}
Exercise 19
\end{exercise}
\begin{hint}
% VIII-PEDAGOGY-HINT-v1.1 exr:viii23-19 replace_template
For a chain isomorphism \(u:C\to D\), conjugate the contraction to \(usu^{-1}\) and verify the identity on \(D\).
\end{hint}
\begin{solution}
Contractibility is preserved by chain isomorphism.
\end{solution}

\begin{exercise}\label{exr:viii23-20}
Exercise 20
\end{exercise}
\begin{hint}
% VIII-PEDAGOGY-HINT-v1.1 exr:viii23-20 replace_template
The only maps between a complex and zero compose to zero; the homotopy-inverse condition therefore says its identity is null-homotopic.
\end{hint}
\begin{solution}
A chain-homotopy-equivalent complex to zero is contractible.
\end{solution}

\begin{exercise}\label{exr:viii23-21}
Exercise 21
\end{exercise}
\begin{hint}
% VIII-PEDAGOGY-HINT-v1.1 exr:viii23-21 replace_template
A topological contraction first compares the identity chain map with a map through a point; pass to the reduced or augmented model to eliminate the point's degree-zero class.
\end{hint}
\begin{solution}
Reduced chains align chain contraction with topological contraction.
\end{solution}

\begin{exercise}\label{exr:viii23-22}
Exercise 22
\end{exercise}
\begin{hint}
% VIII-PEDAGOGY-HINT-v1.1 exr:viii23-22 replace_template
Check that \(s\) preserves each filtration level; an unfiltered contraction need not induce a valid contraction of the filtered objects.
\end{hint}
\begin{solution}
Filtration-compatible contractions can simplify filtered complexes.
\end{solution}

\begin{exercise}\label{exr:viii23-23}
Exercise 23
\end{exercise}
\begin{hint}
% VIII-PEDAGOGY-HINT-v1.1 exr:viii23-23 replace_template
For a contraction on the second tensor factor, test \(S(c\otimes d)=(-1)^{|c|}c\otimes s(d)\) and verify cancellation with the tensor differential.
\end{hint}
\begin{solution}
Tensor contractions require grading signs.
\end{solution}

\begin{exercise}\label{exr:viii23-24}
Exercise 24
\end{exercise}
\begin{hint}
% VIII-PEDAGOGY-HINT-v1.1 exr:viii23-24 replace_template
Compare the two tests: an acyclic cone detects a quasi-isomorphism, while an explicit cone contraction detects the stronger chain-homotopy-equivalence condition.
\end{hint}
\begin{solution}
Mapping cones provide the next major test for chain equivalences.
\end{solution}


\section{Solved problem dossiers}

\begin{problem}[Cycle filling]\label{prob:viii23-01}
Show a contraction fills every cycle.
\end{problem}
\begin{solution}
If \(\partial z=0\), then \(z=\partial s(z)+s\partial z=\partial s(z)\).
\end{solution}

\begin{problem}[Elementary pair]\label{prob:viii23-02}
Construct a contraction on \(0\to A\xrightarrow{\mathrm{id}}A\to0\).
\end{problem}
\begin{solution}
Take \(s\) from the lower copy of \(A\) back to the upper copy as the identity and all other components zero.
\end{solution}

\begin{problem}[Split exact]\label{prob:viii23-03}
Explain how splittings yield a contraction.
\end{problem}
\begin{solution}
Decompose each \(C_n\) into a boundary part and a chosen complement mapping isomorphically onto \(B_{n-1}\); invert that isomorphism to define \(s\).
\end{solution}

\begin{problem}[Vector spaces]\label{prob:viii23-04}
Why does acyclic imply contractible over a field?
\end{problem}
\begin{solution}
All short exact sequences of vector spaces split, so the split-exact theorem applies.
\end{solution}

\begin{problem}[Torsion obstruction]\label{prob:viii23-05}
Why is \(0\to\mathbb Z\xrightarrow{\times2}\mathbb Z\to0\) not contractible?
\end{problem}
\begin{solution}
Its \(H_0\) is \(\mathbb Z/2\), but contractible complexes are acyclic.
\end{solution}

\begin{problem}[Simplex contraction]\label{prob:viii23-06}
Describe the vertex-coning contraction on reduced simplicial chains.
\end{problem}
\begin{solution}
Choose a vertex \(v\) and prepend it to each oriented simplex, with the usual sign/convention; the boundary identity gives \(\partial s+s\partial=\mathrm{id}\).
\end{solution}

\begin{problem}[Reduced chains]\label{prob:viii23-07}
Why can the unreduced chain complex of a point not be contracted?
\end{problem}
\begin{solution}
It has \(H_0\cong\mathbb Z\ne0\), contradicting contractible-implies-acyclic.
\end{solution}

\begin{problem}[Map out]\label{prob:viii23-08}
If \(C\) contracts, prove every chain map \(f:C\to D\) is null-homotopic.
\end{problem}
\begin{solution}
Compose the contraction with \(f\): \(f=f\partial s+fs\partial=\partial(fs)+(fs)\partial\).
\end{solution}

\begin{problem}[Map in]\label{prob:viii23-09}
If \(C\) contracts, prove every chain map \(g:D\to C\) is null-homotopic.
\end{problem}
\begin{solution}
Use \(sg\): \(g=\partial(sg)+(sg)\partial\).
\end{solution}

\begin{problem}[Direct sum]\label{prob:viii23-10}
Construct a contraction on \(C\oplus D\) from contractions on each.
\end{problem}
\begin{solution}
Use \(s_C\oplus s_D\).
\end{solution}

\begin{problem}[Unit cancellation]\label{prob:viii23-11}
Why does a unit incidence coefficient permit cancellation?
\end{problem}
\begin{solution}
The corresponding differential component is an isomorphism on rank-one summands, producing an elementary contractible two-term summand.
\end{solution}

\begin{problem}[Field elimination]\label{prob:viii23-12}
Explain finite-dimensional acyclic complexes over a field as sums of elementary pairs.
\end{problem}
\begin{solution}
Choose bases adapted to images and complements; exactness pairs each complement with the next boundary space by an isomorphism.
\end{solution}

\begin{problem}[Smith form]\label{prob:viii23-13}
What does a Smith diagonal entry \(1\) mean homologically?
\end{problem}
\begin{solution}
It yields an isomorphism between rank-one summands, hence no homology and a contractible pair.
\end{solution}

\begin{problem}[Smith torsion]\label{prob:viii23-14}
What does a Smith diagonal entry \(m>1\) mean?
\end{problem}
\begin{solution}
It contributes a cokernel \(\mathbb Z/m\), so the summand is not contractible.
\end{solution}

\begin{problem}[Explicit exact example]\label{prob:viii23-15}
Verify the displayed \(\mathbb Z\to\mathbb Z^2\to\mathbb Z\) complex is exact.
\end{problem}
\begin{solution}
The first image is \(\mathbb Z(1,0)\), exactly the kernel of \((0\ 1)\), and the second map is surjective.
\end{solution}

\begin{problem}[Contraction strength]\label{prob:viii23-16}
Why is a contraction stronger data than vanishing homology?
\end{problem}
\begin{solution}
It gives an explicit degree-one operator solving \(\mathrm{id}=\partial s+s\partial\), not just an existence statement about boundaries.
\end{solution}

\begin{problem}[Homotopy to zero]\label{prob:viii23-17}
Restate chain contractibility in VIII/22 language.
\end{problem}
\begin{solution}
The identity chain map is chain homotopic to the zero chain map.
\end{solution}

\begin{problem}[Invariance]\label{prob:viii23-18}
If \(C\cong D\) as chain complexes and \(C\) contracts, why does \(D\) contract?
\end{problem}
\begin{solution}
Transport the contraction through the chain isomorphism and its inverse.
\end{solution}

\begin{problem}[Topological analogy]\label{prob:viii23-19}
What topological object suggests the chain contraction of a simplex?
\end{problem}
\begin{solution}
A simplex is contractible to a vertex; the coning operator is the chain-level analogue.
\end{solution}

\begin{problem}[Filtration]\label{prob:viii23-20}
Why might a filtration-compatible contraction be useful?
\end{problem}
\begin{solution}
It removes contractible pieces without mixing filtration levels, preserving the structure needed for later filtered calculations.
\end{solution}

\begin{problem}[Bridge to cones]\label{prob:viii23-21}
What should happen to the mapping cone of a chain-homotopy equivalence?
\end{problem}
\begin{solution}
It should be contractible; VIII/24 proves the cone packages failure of invertibility up to chain homotopy.
\end{solution}

\begin{problem}[Acyclic cone preview]\label{prob:viii23-22}
What should happen to the cone of a quasi-isomorphism?
\end{problem}
\begin{solution}
Its homology vanishes, so the cone is acyclic.
\end{solution}

\begin{problem}[Contraction versus quasi-isomorphism]\label{prob:viii23-23}
Why is contractible stronger than quasi-isomorphic to zero?
\end{problem}
\begin{solution}
Quasi-isomorphic to zero only says homology vanishes; contractible supplies a homotopy-level inverse to zero.
\end{solution}

\begin{problem}[Final principle]\label{prob:viii23-24}
State the practical use of a contraction.
\end{problem}
\begin{solution}
It gives explicit algebraic cancellation and a formula for filling cycles.
\end{solution}
